% ===========================================================================
%  comune/frontespizio-sviluppatore.tex
% ===========================================================================
\begin{titlepage}
\thispagestyle{empty}
\begin{tikzpicture}[remember picture, overlay]

  % La banda verde acqua in alto: è il segno del prodotto, e serve a
  % riconoscere il documento di lato in una pila di fogli.
  \fill[lyraAccento] (current page.north west)
    rectangle ([yshift = -9mm]current page.north east);

  \node[anchor = north west, inner sep = 0pt]
    at ([shift = {(28mm, -42mm)}]current page.north west)
    {\begin{minipage}{125mm}
       {\sffamily\fontsize{10}{12}\selectfont
        \color{lyraAccentoScuro}\bfseries
        LYRA CAPITAL\par}
       \vspace{3mm}
       {\sffamily\fontsize{40}{44}\selectfont\color{lyraInchiostro}
        Guida dello\\sviluppatore\par}
       \vspace{5mm}
       {\color{lyraAccento}\rule{34mm}{1.6pt}\par}
       \vspace{6mm}
       {\sffamily\fontsize{12}{17}\selectfont\color{lyraInchiostro!62}
        Architettura, motore quantitativo,\\
        estensioni e rilascio della piattaforma.\par}
     \end{minipage}};

  % Il riquadro in basso dice a chi serve il documento. Sta in copertina
  % perché chi lo apre per sbaglio deve accorgersene subito: il manuale utente
  % (la User guide, in inglese) è un altro file, e cercare qui come si preme
  % un pulsante è tempo perso.
  \node[anchor = south west, inner sep = 0pt]
    at ([shift = {(28mm, 34mm)}]current page.south west)
    {\begin{minipage}{125mm}
       \begin{tcolorbox}[enhanced, colback = lyraPannello,
                         colframe = lyraBordo, boxrule = 0.6pt, arc = 4pt,
                         left = 10pt, right = 10pt, top = 9pt, bottom = 9pt]
         {\sffamily\footnotesize\color{lyraInchiostro}%
          \textbf{A chi serve.} A chi deve \emph{modificare} Lyra: aggiungere
          una strategia, cambiare il motore, toccare l'interfaccia, mettere
          mano al marketplace.\par
          \vspace{4pt}
          \textbf{A chi non serve.} A chi deve \emph{usare} Lyra: quello è la
          \emph{User guide}, il manuale utente in inglese, documento separato
          nella stessa cartella.\par}
       \end{tcolorbox}
     \end{minipage}};

  \node[anchor = south east, inner sep = 0pt, align = right]
    at ([shift = {(-24mm, 18mm)}]current page.south east)
    {\sffamily\scriptsize\color{lyraInchiostro!45}
     Documento generato da sorgente \LaTeX\\
     \today};

\end{tikzpicture}
\end{titlepage}

\cleardoublepage

% --- la pagina che dice come è fatto il documento ---------------------------
\thispagestyle{empty}
\vspace*{18mm}

{\sffamily\Large\bfseries\color{lyraInchiostro}Come è fatto questo documento\par}
\vspace{3mm}
{\noindent\color{lyraAccento}\rule{\textwidth}{1pt}\par}
\vspace{7mm}

\noindent
Quattro parti, in ordine di quanto a fondo si deve andare.

\begin{description}[leftmargin = 2.4em, style = nextline, itemsep = 5pt]
  \item[Parte I — Orientarsi]
    Che cos'è Lyra, come è diviso, e come si mette in moto una copia sulla
    propria macchina. Si legge una volta.
  \item[Parte II — Come è fatto]
    L'architettura e il motore di calcolo: dove stanno le decisioni e perché
    sono dove sono. È la parte da leggere prima di toccare qualunque cosa.
  \item[Parte III — Come si cambia]
    Aggiungere una strategia, una modalità di pesatura, una pagina,
    un'operazione sul database. Ogni capitolo è una procedura che si segue
    con il codice aperto accanto.
  \item[Parte IV — Tenerlo in piedi]
    Le prove automatiche, il rilascio, e l'elenco di ciò che si rompe più
    spesso con il motivo accanto.
\end{description}

\vspace{6mm}
{\sffamily\large\bfseries\color{lyraInchiostro}Le convenzioni\par}
\vspace{3mm}

\noindent
I riquadri sono quattro e hanno quattro significati distinti.

\vspace{3mm}
\begin{nota}{Nota}
Un chiarimento utile ma non necessario: saltandolo non ci si perde niente di
ciò che serve per lavorare.
\end{nota}

\begin{perche}{Perché così}
La ragione di una scelta che, vista da fuori, sembrerebbe potersi fare
diversamente. Serve a non «semplificare» qualcosa che è complicato per un
motivo.
\end{perche}

\begin{pratica}{In pratica}
Cosa fare, passo per passo. Se si cerca solo l'operazione, sono questi i
riquadri da seguire.
\end{pratica}

\begin{attenzione}{Attenzione}
Qui si rompe qualcosa: dati persi, denaro movimentato, o un errore che si
scopre solo settimane dopo. Non sono avvisi di cortesia.
\end{attenzione}

\vspace{4mm}
\noindent
Nel testo, \cod{questo} è un simbolo del codice, \file{questo/} un percorso nel
progetto, \voce{Questa} una voce dell'interfaccia e \pulsante{Questo} un
pulsante da premere. Ogni blocco di codice porta in alto a destra il file da
cui viene: la documentazione si legge con il sorgente aperto accanto, e un
frammento senza il suo file costringe a cercarlo.

\cleardoublepage
